%%%PAGE 131 rot=0 legibility=fair summary=天体物理复习提纲：比值类、体系分类、万有引力、关键词(公转/表面g/表面转：近地卫星、完全失重、赤道与两极)；夹有地磁场小注与零星涂划
%%%BLOCK id=131-01 ch=05 sec=天体物理·复习提纲 kind=summary alt=none cont=none y=0.04-0.27
\textbf{天体物理}

比值类\quad \fbox{先表达，后正比，常量都当1}\quad 注意 $4{:}3$ 和 $3{:}4$ \underline{反应\ 一下}\quad （其上方小字：先看清后下笔）

\textbf{一、体系分类}

1.\ 题型：
\begin{tabular}[t]{@{}c@{\ ；}c@{\ ；}c@{\ ；}c@{\ ；}c@{}}
比值(比大小) & 轨道 & 多星 & 能量 & 引力\\
90\% & 5\% & 1\% & 0.1\% & 4.9\%
\end{tabular}
% 百分数分别写在 比值、轨道、多星、能量、引力 的下方

2.\ 旋转；表面
%%%END
%%%BLOCK id=131-02 ch=99 sec=化学·气体密度 kind=note alt=none cont=none y=0.03-0.06
% 页面右上角，压在页眉线上，绝大部分被页边截断，仅可见末尾（与第127、129页同一位置的同一行字）
\unclear{$\ldots$}\,$\unit{g\cdot L^{-1}}$
%%%END
%%%BLOCK id=131-03 ch=05 sec=万有引力定律 kind=knowledge alt=none cont=none y=0.27-0.43
\textbf{二、万有引力}

1.\ $F=\dfrac{GMm}{r^{2}}$\qquad $G_{\text{\unclear{大}}}$：$6.67\times10^{-11}\unit{N\cdot m^{2}/kg^{2}}$\quad ——\ 铅球；扭秤

\hspace{6em}卡文迪许\qquad\qquad 放大形变
% 右侧图：扭秤示意——上方一根竖线下挂一个小圆，其下另有一个小圆连一条水平横线（右端有一短竖线）
%FIG p131_a 0.815 0.258 0.95 0.358
\notefig[0.25]{figures/p131_a.png}

2.\ $\left\{\begin{array}{l} r\to\infty\ \ F\to0\\ r\to0\ \ F\to\infty\ \ (\times)\ \text{质点} \end{array}\right.$
%%%END
%%%BLOCK id=131-04 ch=05 sec=卫星运行规律 kind=formula alt=none cont=none y=0.43-0.62
\textbf{三、：关键词}

1.\ “\unclear{公}转”
% 图：日(圆圈)与地(圆圈)用线相连，线上标 r；另有一段弧线(轨道)穿过“地”
%FIG p131_b 0.25 0.476 0.385 0.62
\notefig[0.25]{figures/p131_b.png}
% 右侧一个手绘大圈(云状)圈住下面这行，圈内上方写“认为圆周”
\begin{keybox}[认为圆周]
\[ \dfrac{GMm}{r^{2}}=m\dfrac{v^{2}}{r}=m\omega^{2}r=m\dfrac{4\pi^{2}}{T^{2}}r=ma=m\omega v \]
\end{keybox}
\[ v=\sqrt{\dfrac{GM}{r}}\qquad \omega=\sqrt{\dfrac{GM}{r^{3}}}\qquad T=\sqrt{\dfrac{4\pi^{2}r^{3}}{GM}}\qquad a=\dfrac{GM}{r^{2}} \]
%%%END
%%%BLOCK id=131-05 ch=05 sec=天体表面重力加速度 kind=formula alt=none cont=none y=0.62-0.80
2.\ “表面”\,——\,$g$\ 表面重力加速度\qquad\qquad $g\ \ g$
% 图：圆(地)，圆心处写“地”，半径线标 R；圆外右上有一个小圆圈符号；圆下方右侧有“(a)”样标记（其前有一个难辨的字）
%FIG p131_c 0.155 0.665 0.345 0.80
\notefig[0.25]{figures/p131_c.png}
\[ \dfrac{GMm}{R^{2}}=\fbox{$mg$} \]
此时重力为“效果”
\[ \underline{GM=gR^{2}} \]
$\llcorner$\,黄金代换（给了 $g$ 才能用）
%%%END
%%%BLOCK id=131-06 ch=11 sec=地磁场 kind=note alt=none cont=none y=0.62-0.86
% 页面右侧，周围有大片波浪线涂划；“地磁南北是”下方有一道波浪线；其上方另有几个像“m m”的涂写符号
地磁南北是\ 反向、

不是匀强磁场
%%%END
%%%BLOCK id=131-07 ch=05 sec=近地卫星·完全失重 kind=formula alt=02 cont=none y=0.80-1.0
3.\ \underline{表面转}\,/\,\underline{近地卫星}\,/\,\underline{第一宇宙速度}\,/\,\underline{完全失重}
% 左下图：浅灰色大圆内有一个深色扁椭圆，椭圆右端有一个小圆圈（近地卫星轨道示意）
%FIG p131_d 0.12 0.85 0.305 0.962
\notefig[0.25]{figures/p131_d.png}
\[ \dfrac{GMm}{R}=m\dfrac{v^{2}}{R}=m\omega^{2}R=m\dfrac{4\pi^{2}}{T^{2}}R=ma \]
% CHECK: 第一个分式分母写的是 R（应为 R^2），原样转录
% 右上有两个被圈出的符号（左圈内像“40”，右圈内像“日”），圈旁有大量波浪线涂划
（圈出：\unclear{40}、\unclear{日}）
\[ F_{\text{引}}=F_{\text{重}}+F_{\text{向}} \]
% F_重 下方有向下箭头，F_向 下方有向上箭头
$G=0$ 时 \unclear{万$\ldots$} 完全失重
赤道\quad $\struck{F_{\text{引}}}=F_{\text{\unclear{万}}}=F_{\text{向}}+mg_{\text{赤}}$

两极\quad $F=mg$\quad \unclear{$\dfrac{120}{100}\ \dfrac{220}{370}\ 420$}
% 以下两行被手绘多边形(菱形)圈住，周围有大片涂划
\begin{keybox}
越高 $g$ 越小

越往两极 $g$ 越大
\end{keybox}
%%%END
%%%BLOCK id=131-08 ch=99 sec=涂鸦·无关数字 kind=misc alt=none cont=none y=0.88-0.95
% 圆圈内的数字（每个下方有下划线与小点）：100、50、67、150，其后另有一个被箭头划去的 150；其下有一排涂划线（像连写的“6”）
\unclear{100\ \ 50\ \ 67\ \ 150}
% 另：页面左侧(“W”形线)、右下角(网格状)有大片涂划，未转录
%%%END
%%%BLOCK id=131-09 ch=99 sec=数学·函数 kind=misc alt=none cont=none y=0.94-1.0
% 页面左下角，旁有波浪线涂划
\[ \unclear{f(m)}=\dfrac{m^{2}+2m+\unclear{1}}{2m+3} \]
%%%END
%%%PAGEEND 131
